\clearpage
\section{Triangular bunch in a conducting circular pipe}
\label{sec:triangle-pipe}
This control reproduces the longitudinal-force component of the straight
round-pipe triangular-profile benchmark used in ORBIT~\cite{orbit2006,pyorbit3triangle}.
The new explicit \code{CIRCULAR_UNIFORM} model is compared with the
analytic solution, the earlier \code{CIRCULAR} model and repeated PyORBIT
calculations. This is a frozen-field test, not a tune-shift or tracking test.

The source has fixed circular radius $a=5$ mm, full length $L=10$ m,
$N=10^{14}$ protons and kinetic energy $0.4$ GeV. The grounded pipe radius
is $b=20$ mm. With $s=0$ at the bunch centre,
\begin{align}
 \lambda(s)&=\frac{2Q}{L}\left(1-\frac{2|s|}{L}\right),\quad |s|<L/2,
 & E_s(r,s)&=-\frac{\lambda'(s)}{4\pi\eps\gamma^2}\,g(r),\label{eq:triangle-pipe}\\
 g(r)&=\begin{cases}
 1+2\ln(b/a)-r^2/a^2,&0\le r\le a,\\
 2\ln(b/r),&a<r\le b.
 \end{cases}
 \label{eq:triangle-pointwise-g}
\end{align}
The density vanishes outside its support. On axis, $g=3.772589$ and
$E_*=|E_s|=10.681199$ kV/m; its uniform-disk average is
$\overline g=1/2+2\ln(b/a)$~\cite{baartman}.
Acceptance uses the same 62 on-axis and 82 radial targets at
$1\le|s|\le4$ m, excluding the cusp and ends. No coefficient is fitted.
A proton's equivalent one-metre kick is $\Delta W[\mathrm{eV}]=E_s[\mathrm{V/m}]$.

\reportfigure{triangle-pipe-fixed-fields}{Frozen on-axis fields and signed
difference from the physical analytic solution (black dashed). Gray crosses
retain the pre-change CIRCULAR result; blue circles use CIRCULAR\_UNIFORM.
PyORBIT variants are rerun independently. Residuals use the fixed scale
$E_*$; gray regions are excluded and the gray band marks $\pm1\%$.
The displayed cusp transition has no accuracy claim.}{fig:triangle-pipe-fixed-fields}
\tabfile{triangle-fixed-results}
Both columns are maximum pointwise relative physical field errors; accepted
reference values are nonzero. The fine new OPALX model meets 1\% on both
target sets. Its maximum differences from PyORBIT slices and Baartman
are 0.447\% and 0.167\%, respectively, relative to the local analytic field.

\clearpage
\subsection{Separate physical geometry from the numerical mesh}
The pointwise model requires explicit \code{PIPERADIUS}$=b\ge a$ and
\code{BEAMR}$=a$. \code{PIPESIZEX/Y} specify full open numerical widths
$W_x,W_y$. The physical radius is independent of these widths.
\begin{figure}[H]
\centering
\begin{minipage}[c]{0.53\linewidth}
\small
Model-specific settings for this benchmark:
\begin{verbatim}
FS: FIELDSOLVER, TYPE="FFT2D5",
 PIPEMODE="CIRCULAR_UNIFORM",
 BEAMR=0.005, PIPERADIUS=0.020,
 PIPESIZEX=0.020, PIPESIZEY=0.020,
 NX=128, NY=128, NZ=128,
 PARFFTX=FALSE, PARFFTY=FALSE,
 PARFFTZ=FALSE,
 REFPATHFNAME="straight-path.dat";
\end{verbatim}
\end{minipage}\hfill
\begin{minipage}[c]{0.44\linewidth}
\centering\begin{tikzpicture}[scale=.72,>=Stealth,font=\small]
 \fill[gray!18] (0,0) circle[radius=.5];
 \draw[gray!40,step=.25,very thin] (-1,-1) grid (1,1);
 \draw[blue!65!black,dashed,thick] (-1,-1) rectangle (1,1);
 \draw[black,thick] (0,0) circle[radius=2];
 \draw[gray!75] (0,0) circle[radius=.5];
 \draw[->,gray!65] (-2.2,0)--(2.35,0) node[right] {$x$};
 \draw[->,gray!65] (0,-2.2)--(0,2.35) node[above] {$y$};
 \draw[<->,black] (0,0)--(1.4142,1.4142)
   node[pos=.65,above,sloped,fill=white,inner sep=1pt] {$b=20$ mm};
 \draw[<->,black] (0,0)--(.5,0);
 \node[anchor=west,fill=white,inner sep=1pt] at (.57,.27) {$a=5$ mm};
 \draw[<->,blue!65!black] (-1,-1.15)--(1,-1.15)
   node[midway,below,fill=white,inner sep=1pt] {$W=20$ mm};
 \node[font=\footnotesize,align=center] at (0,-2.65)
   {Physical pipe: circle $b$\\Open numerical mesh: square $W$};
\end{tikzpicture}

\end{minipage}
\caption{The circle is the physical longitudinal-model pipe; the square is
the open transverse numerical grid. Distribution, beam and tracking commands
are supplied by the archived complete input. A transverse conducting wall
is not solved. For this centred round source, the local transverse image
field vanishes.}
\label{fig:pipe-model-geometry}
\end{figure}
\reportfigure{triangle-pipe-fixed-controls}{Left: radial dependence at
$s=+2.5$ m for the same selected results. Right: physical on-axis errors
under matched mesh changes. OPALX uses one source over $64^3,128^3,256^3$;
the historical PyORBIT transverse scan holds $N_z=128$ and its 16-ring source
fixed. The primary PyORBIT curve on the left uses 32 rings. Zero remains
visible; no monotonic or joint asymptotic convergence is claimed.}
{fig:triangle-pipe-fixed-controls}

All numerical changes quoted here use the 144 accepted targets. The last
matched OPALX mesh change is 0.1373\% of $E_*$. Doubling numerical width
and transverse cell count at fixed $b$ changes the field by less than
$5.5\times10^{-12}\%$; changing physical $b$ from 20 to 10 mm instead
reproduces its distinct analytic field. The initial PyORBIT ring-quadrature
bias and longitudinal sampling controls remain archived; repeated primary
results reproduce their earlier values.

Legacy \code{CIRCULAR} retains the moment coefficient
$0.67+2\ln(b/a)$~\cite{oeftiger}, with its historical width fallback when
\code{PIPERADIUS} is absent. Its unchanged lack of radial dependence explains
the earlier pointwise discrepancy. A matched legacy rerun with the new
executable agrees with its archived result within $7.4\times10^{-13}\%$ of
$E_*$. The pipe model is selected explicitly and retains the legacy CIRCULAR
closure. The later OPEN replacement is tested separately in \cref{sec:free-open}.

\clearpage
\subsection{Independent unit tests and an OPALX regression run}
\label{sec:pipe-gates}
The new physics unit tests compute the reference by integrating Gauss' law
to the grounded wall and differentiating the potential longitudinally.
They cover inside-beam, vacuum-annulus and wall fields, a non-unit Lorentz
factor, physical-radius versus mesh-width changes and rejection of particles
outside the physical pipe. Configuration tests check the explicit radius
contract and reject unsupported or inconsistent settings.

\begin{center}\small
\begin{tabular}{>{\raggedright\arraybackslash}p{0.29\linewidth}>{\raggedright\arraybackslash}p{0.63\linewidth}}
\toprule Check & Result and scope\\\midrule
C++ unit tests & 100 pass; one MPI-only test is skipped in the single-rank run.
The addition contains three independent physics tests and seven configuration
tests.\\
Reference quadrature & Independent integration/differentiation check:
maximum relative discrepancy $1.06\times10^{-11}$, below the justified
$2\times10^{-8}$ unit tolerance.\\
Native study & Six new OPALX source/probe pairs and three repeated PyORBIT
configurations. Fields, charge, positions, input and binary hashes are checked.\\
Actual OPALX regression & \code{FFT2D5-TriangularPipe} executes OPALX and checks
its H5 fields at eight nonzero targets plus two wall targets against an
independent Gauss-integral reference. Maximum physical error is 0.00352\%;
wall field is exactly zero.\\
Negative control & The same native case with the legacy model is rejected
by the physical validator; a probe-only solve confirms negligible longitudinal
probe contribution.\\
Smoke integration & STAT reference checks and the field-physics validator pass
through the local smoke runner. No remote CI execution is claimed.\\
Historical OPEN compatibility & At the pipe-update executable, ZeroSC and 1 pC regressions pass all
128 samples and six STAT reference checks with the new binary. The largest
RMS difference from their stored baseline is $4.05\times10^{-10}$ relative.\\
\bottomrule
\end{tabular}
\end{center}
The native field tolerance remains 1\%; wall targets use an absolute
$10^{-4}$ V/m tolerance because their analytic field is zero.
OPALX uses a $10^{-20}$ s readout and matched probe subtraction in the full
study. PyORBIT uses zero-charge probes and fixed-position kicks.
The repeated PyORBIT one-/two-rank fields agree within
$2.6\times10^{-12}\%$ of $E_*$; FFT2D5 remains single-rank.

The earlier benchmark is identified by binary
\begin{quote}\small
\path{e85bc09135bef40c6e0b9a278f45d9aab05ad595f2635e7c2088526bb1de70bf};
\end{quote}
the explicit pointwise model and its native tests use
\begin{quote}\small
\path{9d3771d58f0de7bf913909955e4ef81aef69a33341202edefa9553bb78ead8b5}.
\end{quote}
The before/after fine comparison uses the same physical source, targets and
$256^3$ mesh. Previous tracking and curved-tube results retain their original
binaries and interpretations.

PyORBIT slices obtain longitudinal kicks by differentiating the same grounded
slice potential used for their transverse kicks. OPALX instead combines the
analytic uniform-round $g(r)$ with its sampled $\lambda'$ and open transverse
Poisson solves. Their continuum fields coincide for this prescribed source;
this comparison does not establish a general common-potential PIC map or
symplectic tracking.

\textbf{Limit.} This validates a fixed-radius, coaxial uniform round beam's
local longitudinal pipe model on its linear density branches. Transverse
slices remain open. The model does not add nonlocal end fields, a curvature
operator or a general conducting-pipe solution. Those historical OPEN reruns test backward
compatibility only; full-drift and curved-source accuracy of the new
\code{CIRCULAR_UNIFORM} model are not established.
